\chapter{2011年新课标卷（理）}

\section{单选题}

\begin{problem}
复数 \(\frac{2 + \i}{1 - 2\i}\) 的共轭复数是（\quad）

\choices
    {\(-\frac{3}{5}\i\)}
    {\(\frac{3}{5}\i\)}
    {\(-\i\)}
    {\(\i\)}
\end{problem}

\begin{answer}
C.
\end{answer}

\begin{solution}
将分式的分子、分母同乘以分母的共轭复数，得
\[
\frac{2+\i}{1-2\i}=\frac{(2+\i)(1+2\i)}{(1-2\i)(1+2\i)}
=\frac{5\i}{5}=\i
\]
因此其共轭复数为 \(-\i\)，选 C.
\end{solution}

\begin{problem}
下列函数中，既是偶函数又在 \(\left(0, +\infty\right)\) 单调递增的函数是（\quad）

\choices
    {\(y = x^{3}\)}
    {\(y = |x| + 1\)}
    {\(y = -x^{2} + 1\)}
    {\(y = 2^{-|x|}\)}
\end{problem}

\begin{answer}
B.
\end{answer}

\begin{solution}
逐项判断：\(x^{3}\) 是奇函数；\(|x|+1\) 是偶函数，且在 \(\left(0,+\infty\right)\) 上随 \(x\) 增大而增大；\(-x^{2}+1\) 在该区间单调递减；\(2^{-|x|}\) 虽为偶函数，但在该区间单调递减. 因此选 B.
\end{solution}

\begin{problem}
执行如图的程序框图，如果输入的 \(N\) 是 \(6\)，那么输出的 \(p\) 是（\quad）

\bitmapfigure[max width=0.46\linewidth]{img/2011/new_curriculum_science/q3_fig1.png}

\choices
    {\(120\)}
    {\(720\)}
    {\(1440\)}
    {\(5040\)}
\end{problem}

\begin{answer}
B.
\end{answer}

\begin{solution}
程序先令 \(k=1\)，\(p=1\). 每次先执行 \(p\leftarrow pk\)，再判断 \(k<N\)；若判断为“是”，才执行 \(k\leftarrow k+1\)，然后返回执行乘法运算. 当 \(N=6\) 时，各次执行乘法后的状态依次为
\((k,p)=(1,1)\)，\((2,2)\)，\((3,6)\)，\((4,24)\)，\((5,120)\)，\((6,720)\).

当 \(k=6\) 时不满足 \(k<N\)，程序输出
\[
p=720
\]
因此选 B.
\end{solution}

\begin{problem}
有 \(3\) 个兴趣小组，甲，乙两位同学各自参加其中一个小组，每位同学参加各个小组的可能性相同，则这两位同学参加同一个兴趣小组的概率为（\quad）

\choices
    {\(\frac{1}{3}\)}
    {\(\frac{1}{2}\)}
    {\(\frac{2}{3}\)}
    {\(\frac{3}{4}\)}
\end{problem}

\begin{answer}
A.
\end{answer}

\begin{solution}
甲、乙分别选择一个小组，所有有序结果共有 \(3\times3=9\) 种，等可能的同组结果为 \((1,1)\)，\((2,2)\)，\((3,3)\)，共有 \(3\) 种. 因此所求概率为
\[
\frac{3}{9}=\frac{1}{3}
\]
选 A.
\end{solution}

\begin{problem}
已知角 \(\theta\) 的顶点与原点重合，始边与 \(x\) 轴的正半轴重合，终边在直线 \(y = 2x\) 上，则 \(\cos 2\theta =\)（\quad）

\choices
    {\(-\frac{4}{5}\)}
    {\(-\frac{3}{5}\)}
    {\(\frac{3}{5}\)}
    {\(\frac{4}{5}\)}
\end{problem}

\begin{answer}
B.
\end{answer}

\begin{solution}
终边在直线 \(y=2x\) 上，所以 \(\tan\theta=2\). 利用二倍角公式，得
\[
\cos2\theta=\frac{1-\tan^{2}\theta}{1+\tan^{2}\theta}
=\frac{1-4}{1+4}=-\frac35
\]
因此选 B.
\end{solution}

\begin{problem}
在一个几何体的三视图中，正视图和俯视图如下图所示，则相应的侧视图可以为（\quad）

\bitmapfigure[max width=0.42\linewidth]{img/2011/new_curriculum_science/q6_condition.png}

\choices
    {\choicebitmap[width=0.15\paperwidth]{img/2011/new_curriculum_science/q6_optA.png}}
    {\choicebitmap[width=0.15\paperwidth]{img/2011/new_curriculum_science/q6_optB.png}}
    {\choicebitmap[width=0.15\paperwidth]{img/2011/new_curriculum_science/q6_optC.png}}
    {\choicebitmap[width=0.15\paperwidth]{img/2011/new_curriculum_science/q6_optD.png}}
\end{problem}

\begin{answer}
D.
\end{answer}

\begin{solution}
由正视图与俯视图可知，几何体的外轮廓在侧视方向应为等腰三角形；俯视图中的中间棱位于对称轴上，且该棱在侧视图中应为可见的实线. A、B 的外轮廓为梯形，C 的外轮廓虽为三角形但中间棱画成虚线，只有 D 同时满足外轮廓、对称轴位置以及可见性要求，因此选 D.
\end{solution}

\begin{problem}
设直线 \(l\) 过双曲线 \(C\) 的一个焦点，且与 \(C\) 的一条对称轴垂直，\(l\) 与 \(C\) 交于 \(A\)，\(B\) 两点，\(|AB|\) 为 \(C\) 的实轴长的 \(2\) 倍，则 \(C\) 的离心率为（\quad）

\choices
    {\(\sqrt{2}\)}
    {\(\sqrt{3}\)}
    {\(2\)}
    {\(3\)}
\end{problem}

\begin{answer}
B.
\end{answer}

\begin{solution}
设双曲线方程为 \(\dfrac{x^{2}}{a^{2}}-\dfrac{y^{2}}{b^{2}}=1\)，焦距参数为 \(c\)，则 \(c^{2}=a^{2}+b^{2}\). 若过焦点的直线垂直于虚轴，则它是实轴方向的直线；与双曲线的交弦只能是实轴，长度为实轴长 \(2a\)，不可能等于 \(4a\). 故该直线垂直于实轴，即为 \(x=c\).

代入双曲线方程，交点的纵坐标为 \(y=\pm b^{2}/a\)，所以 \(|AB|=\frac{2b^{2}}{a}=4a\)，\(b^{2}=2a^{2}\).
从而 \(c^{2}=3a^{2}\)，离心率
\[
e=\frac ca=\sqrt3
\]
因此选 B.
\end{solution}

\begin{problem}
\(\left(x + \frac{a}{x}\right)\left(2x - \frac{1}{x}\right)^5\) 的展开式中各项系数的和为 \(2\)，则该展开式中常数项为（\quad）

\choices
    {\(-40\)}
    {\(-20\)}
    {\(20\)}
    {\(40\)}
\end{problem}

\begin{answer}
D.
\end{answer}

\begin{solution}
展开式各项系数的和等于令 \(x=1\) 后的值，因此
\[
(1+a)(2-1)^{5}=2
\]
得 \(a=1\). 于是
\[
\left(x+\frac1x\right)\left(2x-\frac1x\right)^{5}
\]
常数项有两种来源：从第一个因子的 \(x\) 取第二个因子中 \(x^{-1}\) 项，需要取三次 \(-x^{-1}\)，贡献
\[
\mathrm C_5^3 2^{2}(-1)^{3}=-40
\]
从第一个因子的 \(x^{-1}\) 取第二个因子中 \(x\) 项，需要取两次 \(-x^{-1}\)，贡献
\[
\mathrm C_5^2 2^{3}=80
\]
常数项为 \(-40+80=40\)，因此选 D.
\end{solution}

\begin{problem}
由曲线 \(y = \sqrt{x}\)，直线 \(y = x - 2\) 及 \(y\) 轴所围成的图形的面积为（\quad）

\choices
    {\(\frac{10}{3}\)}
    {\(4\)}
    {\( \frac{16}{3} \)}
    {\(6\)}
\end{problem}

\begin{answer}
C.
\end{answer}

\begin{solution}
令 \(t=\sqrt{x}\)，则交点满足 \(t=t^{2}-2\)，故 \(t=2\)，即 \(x=4\). 在 \(0\leqslant x\leqslant4\) 上，曲线 \(y=\sqrt{x}\) 在直线 \(y=x-2\) 上方，所以面积为
\[
\int_{0}^{4}\left(\sqrt{x}-(x-2)\right)\,\mathrm{d}x
=\left[\frac23x^{3/2}-\frac{x^{2}}2+2x\right]_{0}^{4}
=\frac{16}{3}
\]
因此选 C.
\end{solution}

\begin{problem}
已知 \(\bs{a}\) 与 \(\bs{b}\) 均为单位向量，其夹角为 \(\theta\)，有下列四个命题：
\begin{circlelist}
    \item \(p_1: |\bs{a} + \bs{b}| > 1 \Leftrightarrow \theta \in \left[0, \frac{2\pi}{3}\right)\)
    \item \(p_2: |\bs{a} + \bs{b}| > 1 \Leftrightarrow \theta \in \left(\frac{2\pi}{3}, \pi\right]\)
    \item \(p_3: |\bs{a} - \bs{b}| > 1 \Leftrightarrow \theta \in \left[0, \frac{\pi}{3}\right)\)
    \item \(p_4: |\bs{a} - \bs{b}| > 1 \Leftrightarrow \theta \in \left(\frac{\pi}{3}, \pi\right]\)
\end{circlelist}
其中的真命题是（\quad）

\choices
    {\(p_1\)，\(p_4\)}
    {\(p_1\)，\(p_3\)}
    {\(p_2\)，\(p_3\)}
    {\(p_2\)，\(p_4\)}
\end{problem}

\begin{answer}
A.
\end{answer}

\begin{solution}
因为 \(\bs{a}\)，\(\bs{b}\) 都是单位向量，且两向量的夹角满足 \(0\leqslant\theta\leqslant\pi\)，所以
\[
|\bs{a}+\bs{b}|^{2}=|\bs{a}|^{2}+|\bs{b}|^{2}+2\bs{a}\cdot\bs{b}=2+2\cos\theta
\]
\[
|\bs{a}-\bs{b}|^{2}=|\bs{a}|^{2}+|\bs{b}|^{2}-2\bs{a}\cdot\bs{b}=2-2\cos\theta
\]
由于两边长度均为非负数，分别比较平方可得
\[
|\bs{a}+\bs{b}|>1\Longleftrightarrow 2+2\cos\theta>1
\Longleftrightarrow \cos\theta>-\frac12
\Longleftrightarrow \theta\in\left[0,\frac{2\pi}{3}\right)
\]
所以 \(p_{1}\) 为真，而 \(p_{2}\) 所给区间与上述等价条件不符，故 \(p_{2}\) 为假. 对差向量，
\[
|\bs{a}-\bs{b}|>1\Longleftrightarrow 2-2\cos\theta>1
\Longleftrightarrow \cos\theta<\frac12
\Longleftrightarrow \theta\in\left(\frac{\pi}{3},\pi\right]
\]
所以 \(p_{4}\) 为真，而 \(p_{3}\) 所给区间与上述等价条件相反，故 \(p_{3}\) 为假. 因此真命题为 \(p_{1}\)，\(p_{4}\)，选 A.
\end{solution}

\begin{problem}
设函数 \( f(x) = \sin\left(\omega x + \varphi\right) + \cos\left(\omega x + \varphi\right) \)（\(\omega > 0\)，\(|\varphi| < \frac{\pi}{2}\)）的最小正周期为 \( \pi \)，且 \( f(-x) = f(x) \)，则（\quad）

\choices
    {\(f(x)\) 在 \(\left(0, \frac{\pi}{2}\right)\) 单调递减}
    {\(f(x)\) 在 \(\left(\frac{\pi}{4}, \frac{3\pi}{4}\right)\) 单调递减}
    {\(f(x)\) 在 \(\left(0, \frac{\pi}{2}\right)\) 单调递增}
    {\(f(x)\) 在 \(\left(\frac{\pi}{4}, \frac{3\pi}{4}\right)\) 单调递增}
\end{problem}

\begin{answer}
A.
\end{answer}

\begin{solution}
将函数化为
\[
f(x)=\sqrt2\sin\left(\omega x+\varphi+\frac\pi4\right)
\]
其最小正周期为 \(\pi\)，故 \(\omega=2\). 令 \(\eta=\varphi+\frac\pi4\)，偶性要求
\[
\sin(-2x+\eta)=\sin(2x+\eta)
\]
比较展开式中的 \(\sin2x\)、\(\cos2x\) 系数，得 \(\cos\eta=0\). 结合 \(|\varphi|<\frac\pi2\)，只能取 \(\eta=\frac\pi2\)，从而 \(\varphi=\frac\pi4\)，且
\[
f(x)=\sqrt2\cos2x
\]
在 \(\left(0,\frac\pi2\right)\) 上，\(f'(x)=-2\sqrt2\sin2x<0\)，故选 A.
\end{solution}

\begin{problem}
函数 \( y = \frac{1}{1 - x} \) 的图象与函数 \( y = 2\sin \pi x \)（\(-2 \leqslant x \leqslant 4\)）的图象所有交点的横坐标之和等于（\quad）

\choices
    {\(2\)}
    {\(4\)}
    {\(6\)}
    {\(8\)}
\end{problem}

\begin{answer}
D.
\end{answer}

\begin{solution}
交点横坐标满足
\((1-x)\sin\pi x=\frac12\)，\(x\ne1\).
记左端为 \(G(x)\). 由 \(\sin(\pi(2-x))=-\sin(\pi x)\)，可得 \(G(2-x)=G(x)\)，所以根关于 \(x=1\) 对称. 先考察四个可能使 \(G(x)>0\) 的区间：
这些区间为 \((-2,-1)\)、\((0,1)\)、\((1,2)\) 和 \((3,4)\).
在每个区间内作平移 \(t\in(0,1)\)，在 \((-2,-1)\) 和 \((0,1)\) 上分别得到 \(g_3(t)=(3-t)\sin\pi t\) 和 \(g_1(t)=(1-t)\sin\pi t\)；在 \((1,2)\) 和 \((3,4)\) 上分别得到 \(h_0(t)=t\sin\pi t=g_1(1-t)\) 和 \(h_2(t)=(2+t)\sin\pi t=g_3(1-t)\).
对 \(g_c(t)=(c-t)\sin\pi t\)（\(c=1\)，\(3\)），在 \(0<t<\frac12\) 时，导数为零等价于
\[
\tan\pi t=\pi(c-t)
\]
左端严格递增，右端严格递减，且导数在 \(t=0\) 附近为正、在 \(t=\frac12\) 处为负，故 \(g_c\) 在 \(\left(0,\frac12\right)\) 内有唯一极大值；在 \(\left(\frac12,1\right)\) 内导数恒为负. 又 \(g_3(\frac12)=\frac52>\frac12\)，而 \(g_1(\frac12)=\frac12\) 且 \(g_1'(\frac12)=-1<0\)，所以这两个函数的极大值都大于 \(\frac12\). 由于它们在 \(t=0\)，\(1\) 处均为 \(0\)，方程 \(g_c(t)=\frac12\) 在每个区间内各有两个根. \(h_0\)，\(h_2\) 分别是 \(g_1\)，\(g_3\) 关于 \(t=\frac12\) 的镜像，因此也各有两个根. 由此四个区间共有 \(8\) 个根.

由根的对称性，每一对对称根的和为 \(2\)，共有四对，故所有交点横坐标之和为 \(4\times2=8\)，选 D.
\end{solution}

\section{填空题}

\begin{problem}
若变量 \(x\)，\(y\) 满足约束条件 \(\begin{cases}
3 \leqslant 2x + y \leqslant 9,\\
6 \leqslant x - y \leqslant 9.
\end{cases}\) 则 \(z = x + 2y\) 的最小值为\fillinblank{}.
\end{problem}

\begin{answer}
\(-6\)
\end{answer}

\begin{solution}
令
\(u=2x+y\)，\(v=x-y\).
则 \(3\leqslant u\leqslant9\)，\(6\leqslant v\leqslant9\)，且由 \(x=(u+v)/3\)、\(y=(u-2v)/3\) 得
\(z=x+2y=u-v\).
所以 \(z\) 的最小值在 \(u=3\)，\(v=9\) 时取得，为 \(3-9=-6\).
\end{solution}

\begin{problem}
在平面直角坐标系 \(xOy\) 中，椭圆 \(C\) 的中心为原点，焦点 \(F_{1}\)，\(F_{2}\) 在 \(x\) 轴上，离心率为 \(\frac{\sqrt{2}}{2}\)，过 \(F_{1}\) 的直线 \(l\) 交 \(C\) 于 \(A\)，\(B\) 两点，且 \(\triangle ABF_{2}\) 的周长为 \(16\)，那么 \(C\) 的方程为\fillinblank{}.
\end{problem}

\begin{answer}
\(\frac{x^{2}}{16} + \frac{y^{2}}{8} = 1\)
\end{answer}

\begin{solution}
设椭圆的长半轴、短半轴、半焦距分别为 \(a\)，\(b\)，\(c\)，则 \(e=c/a=\sqrt2/2\)，从而 \(c^{2}=a^{2}/2\)、\(b^{2}=a^{2}-c^{2}=a^{2}/2\). 由椭圆定义，\(AF_{1}+AF_{2}=2a\)，\(BF_{1}+BF_{2}=2a\). 又因 \(F_{1}\) 在线段 \(AB\) 上，\(AB=AF_{1}+BF_{1}\)，故
\[
AB+AF_{2}+BF_{2}=4a=16
\]
得 \(a=4\)，\(b^{2}=8\). 所以椭圆方程为
\[
\frac{x^{2}}{16}+\frac{y^{2}}8=1
\]
\end{solution}

\begin{problem}
已知矩形 \(ABCD\) 的顶点都在半径为 \(4\) 的球 \(O\) 的球面上，且 \(AB = 6\)，\(BC = 2\sqrt{3}\)，则棱锥 \(O - ABCD\) 的体积为\fillinblank{}.
\end{problem}

\begin{answer}
\(8\sqrt{3}\)
\end{answer}

\begin{solution}
矩形四个顶点共圆于其所在平面，矩形中心记为 \(H\). 由于球心 \(O\) 到四个顶点等距，\(O\) 在矩形平面上的投影为其外接圆圆心 \(H\)，故 \(OH\) 垂直于底面. 球心到各顶点的距离为 \(4\)，而
\[
AH=\frac12\sqrt{AB^{2}+BC^{2}}
=\frac12\sqrt{36+12}=2\sqrt3
\]
在直角三角形 \(OHA\) 中，棱锥的高为
\[
OH=\sqrt{OA^{2}-AH^{2}}=\sqrt{16-12}=2
\]
底面面积为 \(6\cdot2\sqrt3=12\sqrt3\)，故体积
\[
V=\frac13\cdot12\sqrt3\cdot2=8\sqrt3
\]
\end{solution}

\begin{problem}
在 \(\triangle ABC\) 中，\(B = 60^{\circ}\)，\(AC = \sqrt{3}\)，则 \(AB + 2BC\) 的最大值为\fillinblank{}.
\end{problem}

\begin{answer}
\(2\sqrt{7}\)
\end{answer}

\begin{solution}
设 \(BC=a\)、\(AB=c\). 由余弦定理及 \(B=60^{\circ}\)、\(AC=\sqrt3\)，有
\[
3=a^{2}+c^{2}-2ac\cos60^{\circ}=a^{2}+c^{2}-ac
\]
利用恒等式
\[
\frac{28}{3}(a^{2}+c^{2}-ac)-(2a+c)^{2}=\frac{(4a-5c)^{2}}3\geqslant0
\]
可得 \(2a+c\leqslant2\sqrt7\). 当 \(4a=5c\) 时取等号，且正数 \(a=5/\sqrt7\)，\(c=4/\sqrt7\) 满足上面的余弦定理条件，故最大值为 \(2\sqrt7\).
\end{solution}

\section{解答题}

\begin{problem}
等比数列 \(\{a_{n}\}\) 的各项均为正数，且 \(2a_{1} + 3a_{2} = 1\)，\(a_{3}^{2} = 9a_{2}a_{6}\).
\begin{enumerate}
\item 求数列 \( \{a_{n}\} \) 的通项公式；
\item 设 \( b_{n} = \log_{3}a_{1} + \log_{3}a_{2} + \cdots + \log_{3}a_{n} \)，求数列 \( \left\{\frac{1}{b_{n}}\right\} \) 的前 \( n \) 项和.
\end{enumerate}
\end{problem}

\begin{solution}
\begin{enumerate}
\item 设等比数列的公比为 \(q>0\). 由 \(a_{3}^{2}=a_{1}^{2}q^{4}\)，\(9a_{2}a_{6}=9a_{1}^{2}q^{6}\)
及各项为正，得 \(q^{2}=1/9\)，从而 \(q=1/3\). 再由
\[
2a_{1}+3a_{2}=a_{1}(2+3q)=1
\]
得 \(a_{1}=1/3\)，所以
\[
a_{n}=a_{1}q^{n-1}=\frac{1}{3^{n}}
\]
\item 由 \(a_{i}=3^{-i}\)，有
\[
b_{n}=\sum_{i=1}^{n}\log_{3}3^{-i}
=-\sum_{i=1}^{n}i=-\frac{n(n+1)}2
\]
因此
\[
\frac1{b_{i}}=-\frac{2}{i(i+1)}
=-2\left(\frac1i-\frac1{i+1}\right)
\]
前 \(n\) 项和为
\[
\sum_{i=1}^{n}\frac1{b_{i}}
=-2\left(1-\frac1{n+1}\right)
=-\frac{2n}{n+1}
\]
\end{enumerate}
\end{solution}

\begin{problem}
如图，四棱锥 \(P - ABCD\) 中，底面 \(ABCD\) 为平行四边形，\(\angle DAB = 60^{\circ}\)，\(AB = 2AD\)，\(PD \perp\) 底面 \(ABCD\).
\begin{enumerate}
\item 证明： \(PA \perp BD\)；
\item 若 \(PD = AD\)，求二面角 \(A - PB - C\) 的余弦值.
\end{enumerate}

\bitmapfigure[max width=0.42\linewidth]{img/2011/new_curriculum_science/q18_fig1.png}
\end{problem}

\begin{solution}
\begin{enumerate}
\item 设 \(AD=a\)，则 \(AB=2a\). 在底面平行四边形中，
\[
BD^{2}=AD^{2}+AB^{2}-2AD\cdot AB\cos60^\circ=3a^{2}
\]
又 \(PD\perp\) 底面，故 \(PD\perp BD\)，从而
\[
AB^{2}=AD^{2}+BD^{2}
\]
得 \(AD\perp BD\). 因为 \(PD\perp\) 底面，故 \(PD\perp AD\)；直线 \(BD\) 同时垂直于平面 \(PAD\) 内相交的 \(AD\)，\(PD\)，所以 \(BD\perp\) 平面 \(PAD\)，进而 \(PA\perp BD\).
\item 取 \(AD=1\)，建立以 \(D\) 为原点、\(DA\)，\(DB\)，\(DP\) 为三个互相垂直坐标轴的直角坐标系. 由 \(BD=\sqrt3\,AD\) 及 \(PD=AD\)，可设 \(A(1,0,0)\)，\(B(0,\sqrt3,0)\)，\(C(-1,\sqrt3,0)\)，\(P(0,0,1)\).
取 \(PB\) 方向向量 \(\bs{w}=(0,\sqrt3,-1)\). 在平面 \(PAB\) 内且垂直于 \(PB\) 的向量可取
\[
\bs{u}=(A-P)-\frac{(A-P)\cdot\bs{w}}{|\bs{w}|^{2}}\bs{w}
=\left(1,-\frac{\sqrt3}{4},-\frac34\right)
\]
在平面 \(PBC\) 内且垂直于 \(PB\) 的向量可取
\[
\bs{v}=(C-P)-\frac{(C-P)\cdot\bs{w}}{|\bs{w}|^{2}}\bs{w}
=(-1,0,0)
\]
这两个向量分别指向两半平面内的 \(A\)，\(C\)，所以它们的夹角就是二面角 \(A-PB-C\). 于是 \(\bs{u}\cdot\bs{v}=-1\)，\(|\bs{u}|=\frac{\sqrt7}{2}\)，\(|\bs{v}|=1\)，
故
\[
\cos\angle(A-PB-C)
=\frac{-1}{(\sqrt7/2)\cdot1}
=-\frac{2\sqrt7}{7}
\]
\end{enumerate}
\end{solution}

\begin{problem}
某种产品的质量以其质量指标值衡量，质量指标值越大表明质量越好，且质量指标值大于或等于 \(102\) 的产品为优质品. 现用两种新配方（分别称为 A 配方和 B 配方）做试验，各生产了 \(100\) 件这种产品，并测量了每件产品的质量指标值，得到了下面试验结果：

\begin{center}
A 配方的频数分布表

\begin{tabular}{c|c|c|c|c|c}
\hline
指标值分组 & \(\left[90,94\right)\) & \(\left[94,98\right)\) & \(\left[98,102\right)\) & \(\left[102,106\right)\) & \(\left[106,110\right]\) \\
\hline
频数 & \(8\) & \(20\) & \(42\) & \(22\) & \(8\) \\
\hline
\end{tabular}

B 配方的频数分布表

\begin{tabular}{c|c|c|c|c|c}
\hline
指标值分组 & \(\left[90,94\right)\) & \(\left[94,98\right)\) & \(\left[98,102\right)\) & \(\left[102,106\right)\) & \(\left[106,110\right]\) \\
\hline
频数 & \(4\) & \(12\) & \(42\) & \(32\) & \(10\) \\
\hline
\end{tabular}
\end{center}

\begin{enumerate}
\item 分别估计用 A 配方，B 配方生产的产品的优质品率；
\item 已知用 B 配方生产一件产品的利润 \(y\)（单位：元）与其质量指标值 \(t\) 的关系式为 \( y = \begin{cases}
-2, & t < 94,\\
2, & 94 \leqslant t < 102,\\
4, & t \geqslant 102.
\end{cases} \) 从用 B 配方生产的产品中任取一件，其利润记为 \(X\)（单位：元），求 \(X\) 的分布列及数学期望. （以实验结果中质量指标值落入各组的频率作为一件产品的质量指标值落入相应组的概率）.
\end{enumerate}
\end{problem}

\begin{solution}
\begin{enumerate}
\item 优质品的指标值满足 \(t\geqslant102\). 因此 A 配方的优质品数为 \(22+8=30\)，B 配方的优质品数为 \(32+10=42\)，相应的优质品率分别为 \(\frac{30}{100}=0.30\)，\(\frac{42}{100}=0.42\).
\item 按 B 配方的频数表及利润分段，随机变量 \(X\) 的分布列为
\begin{center}
\begin{tabular}{c|c|c|c}
\hline
\(X\) & \(-2\) & \(2\) & \(4\)\\
\hline
\(P\) & \(\dfrac4{100}\) & \(\dfrac{12+42}{100}\) & \(\dfrac{32+10}{100}\)\\
\hline
\end{tabular}
\end{center}
即 \(P(X=-2)=0.04\)，\(P(X=2)=0.54\)，\(P(X=4)=0.42\).
所以 \(E(X)=(-2)\times0.04+2\times0.54+4\times0.42=2.68\).
\end{enumerate}
\end{solution}

\begin{problem}
在平面直角坐标系 \(xOy\) 中，已知点 \(A\left(0, -1\right)\)，\(B\) 点在直线 \(y = -3\) 上，\(M\) 点满足 \(\overrightarrow{MB} \parallel \overrightarrow{OA}\)，\(\overrightarrow{MA} \cdot \overrightarrow{AB} = \overrightarrow{MB} \cdot \overrightarrow{BA}\)，\(M\) 点的轨迹为曲线 \(C\).
\begin{enumerate}
\item 求 \(C\) 的方程；
\item \(P\) 为 \(C\) 上的动点，\(l\) 为 \(C\) 在 \(P\) 点处的切线，求 \(O\) 点到 \(l\) 距离的最小值.
\end{enumerate}
\end{problem}

\begin{solution}
\begin{enumerate}
\item 设 \(M(x,y)\). 由 \(\overrightarrow{MB}\parallel\overrightarrow{OA}\) 且 \(B\) 在 \(y=-3\) 上，得 \(B=(x,-3)\). 于是 \(\overrightarrow{MA}=(-x,-1-y)\)，\(\overrightarrow{AB}=(x,-2)\)，\(\overrightarrow{MB}=(0,-3-y)\)，\(\overrightarrow{BA}=(-x,2)\).
代入数量积条件，得
\[
(-x,-1-y)\cdot(x,-2)=(0,-3-y)\cdot(-x,2)
\]
即
\[
-x^{2}+2+2y=-6-2y
\]
所以曲线 \(C\) 的方程为
\[
y=\frac{x^{2}}4-2
\]
\item 设切点为 \(P(x_{0},y_{0})\)，则 \(y_{0}=x_{0}^{2}/4-2\)，切线斜率为 \(x_{0}/2\)，切线方程化为
\[
x_{0}x-2y-\frac{x_{0}^{2}}2-4=0
\]
故原点到该切线的距离为
\[
d=\frac{\frac{x_{0}^{2}}2+4}{\sqrt{x_{0}^{2}+4}}
=\frac12\left(\sqrt{x_{0}^{2}+4}+\frac4{\sqrt{x_{0}^{2}+4}}\right)
\]
令 \(u=\sqrt{x_{0}^{2}+4}\geqslant2\)，由基本不等式
\[
u+\frac4u\geqslant4
\]
得 \(d\geqslant2\)，且 \(x_{0}=0\) 时取等号. 因此最小值为 \(2\).
\end{enumerate}
\end{solution}

\begin{problem}
已知函数 \( f(x) = a \frac{\ln x}{x + 1} + \frac{b}{x} \)，曲线 \( y = f(x) \) 在点 \(\left(1, f(1)\right)\) 处的切线方程为 \( x + 2y - 3 = 0 \).
\begin{enumerate}
    \item 求 \(a\)，\(b\) 的值；
\item 如果当 \(x > 0\)，且 \(x \neq 1\) 时，\(f(x) > \frac{\ln x}{x - 1} + \frac{k}{x}\)，求 \(k\) 的取值范围.
\end{enumerate}
\end{problem}

\begin{solution}
\begin{enumerate}
\item 由切线方程 \(x+2y-3=0\)，得切点纵坐标 \(f(1)=1\)，而
\[
f(1)=b
\]
故 \(b=1\). 又
\[
f'(x)=a\frac{(x+1)/x-\ln x}{(x+1)^{2}}-\frac{b}{x^{2}}
\]
所以 \(f'(1)=a/2-b\). 切线斜率为 \(-1/2\)，从而
\[
\frac a2-1=-\frac12
\]
得 \(a=1\).
\item 代入 \(a=b=1\)，记差值为
\[
D(x)=f(x)-\left(\frac{\ln x}{x-1}+\frac{k}{x}\right)
=\frac{1-k}{x}-\frac{2\ln x}{x^{2}-1}
\]
先取 \(k=0\)，有
\[
D_{0}(x)=\frac{x^{2}-1-2x\ln x}{x(x^{2}-1)}
\]
当 \(x>1\) 时，令 \(\phi(x)=x-x^{-1}-2\ln x\)，则 \(\phi(1)=0\)，\(\phi'(x)=\frac{(x-1)^{2}}{x^{2}}>0\)，
故 \(x^{2}-1-2x\ln x=x\phi(x)>0\)，从而 \(D_{0}(x)>0\). 当 \(0<x<1\) 时，令 \(\psi(x)=1-x^{2}+2x\ln x\)，则
\(\psi(1)=0\)，\(\psi'(x)=2(1-x+\ln x)<0\).
其中用了 \(\ln x<x-1\). 所以 \(x<1\) 时 \(\psi(x)>0\)，即分子 \(x^{2}-1-2x\ln x=-\psi(x)<0\)，分母也小于零，仍有 \(D_{0}(x)>0\).

若 \(k\leqslant0\)，则
\[
D(x)=D_{0}(x)-\frac{k}{x}>0
\]
反之，若 \(k>0\)，则 \(\lim_{x\to1}D(x)=-k<0\)，故在 \(1\) 附近不可能恒有 \(D(x)>0\). 因此
\[
k\in(-\infty,0]
\]
\end{enumerate}
\end{solution}

\begin{problem}
如图，\(D\)，\(E\) 分别为 \(\triangle ABC\) 的边 \(AB\)，\(AC\) 上的点，且不与 \(\triangle ABC\) 的顶点重合. 已知 \(AE\) 的长为 \(m\)，\(AC\) 的长为 \(n\)，\(AD\)，\(AB\) 的长是关于 \(x\) 的方程 \(x^{2} - 14x + mn = 0\) 的两个根.
\begin{enumerate}
\item 证明： \(C\)，\(B\)，\(D\)，\(E\) 四点共圆；
\item 若 \(\angle A = 90^{\circ}\)，且 \(m = 4\)，\(n = 6\)，求 \(C\)，\(B\)，\(D\)，\(E\) 所在圆的半径.
\end{enumerate}

\bitmapfigure[max width=0.38\linewidth]{img/2011/new_curriculum_science/q22_fig1.png}
\end{problem}

\begin{solution}
\begin{enumerate}
\item 因为 \(AD\)，\(AB\) 是方程 \(x^{2}-14x+mn=0\) 的两根，所以
\[
AD\cdot AB=mn=AE\cdot AC
\]
于是
\[
\frac{AD}{AC}=\frac{AE}{AB}
\]
又 \(\angle DAE=\angle CAB\)，故 \(\triangle ADE\sim\triangle ACB\)，从而
\[
\angle ADE=\angle ACB
\]
由于 \(A\)，\(D\)，\(B\) 共线且 \(A\)，\(E\)，\(C\) 共线，\(\angle BDE=180^\circ-\angle ADE\)，\(\angle BCE=\angle ACB\)，所以
\[
\angle BDE+\angle BCE=180^\circ
\]
故 \(C\)，\(B\)，\(D\)，\(E\) 四点共圆.
\item 当 \(m=4\)，\(n=6\) 时，方程为
\[
x^{2}-14x+24=0
\]
其两根为 \(2,12\). 因此 \(AD=2\)，\(AB=12\). 取 \(A(0,0)\)，\(B(12,0)\)，\(C(0,6)\)，\(D(2,0)\)，\(E(0,4)\).
设所求圆方程为 \(x^{2}+y^{2}+ux+vy+w=0\). 代入 \(B\)，\(D\)，\(C\)，分别得
\[
\begin{cases}
144+12u+w=0,\\
4+2u+w=0,\\
36+6v+w=0
\end{cases}
\]
解得 \(u=-14\)，\(v=-10\)，\(w=24\)，所以圆方程为
\[
x^{2}+y^{2}-14x-10y+24=0
\]
或
\[
(x-7)^{2}+(y-5)^{2}=50
\]
故半径为 \(5\sqrt2\).
\end{enumerate}
\end{solution}

\begin{problem}
在直角坐标系 \(xOy\) 中，曲线 \(C_1\) 的参数方程为 \(\begin{cases}
x = 2\cos \alpha,\\
y = 2 + 2\sin \alpha,
\end{cases}\)（\(\alpha\) 为参数），\(M\) 是 \(C_1\) 上的动点，\(P\) 点满足 \(\overrightarrow{OP} = 2\overrightarrow{OM}\)，\(P\) 点的轨迹为曲线 \(C_2\).
\begin{enumerate}
\item 求 \(C_2\) 的方程；
\item 在以 \(O\) 为极点，\(x\) 轴的正半轴为极轴的极坐标系中，射线 \(\theta = \frac{\pi}{3}\) 与 \(C_1\) 的异于极点的交点为 \(A\)，与 \(C_2\) 的异于极点的交点为 \(B\)，求 \(|AB|\).
\end{enumerate}
\end{problem}

\begin{solution}
\begin{enumerate}
\item 由 \(\overrightarrow{OP}=2\overrightarrow{OM}\)，若 \(M=(x_M,y_M)\)、\(P=(x,y)\)，则 \(x=2x_M\)，\(y=2y_M\). 因此
\[
\left(\frac{x}{2}\right)^{2}+\left(\frac{y}{2}-2\right)^{2}=4
\]
即
\[
C_{2}:x^{2}+(y-4)^{2}=16
\]
\item 对 \(C_{1}\) 使用极坐标关系 \(x=r\cos\theta\)，\(y=r\sin\theta\)，得
\[
r^{2}\cos^{2}\theta+(r\sin\theta-2)^{2}=4
\]
整理为
\[
r^{2}-4r\sin\theta=0
\]
即 \(r(r-4\sin\theta)=0\). 异于极点的交点取 \(r\neq0\)，所以在 \(\theta=\pi/3\) 上
\[
r_{A}=4\sin\frac\pi3=2\sqrt3
\]
将 \(C_{2}:x^{2}+(y-4)^{2}=16\) 化为极坐标方程，得
\[
r^{2}\cos^{2}\theta+(r\sin\theta-4)^{2}=16
\]
整理为
\[
r^{2}-8r\sin\theta=0
\]
即 \(r(r-8\sin\theta)=0\). 同样取异于极点的交点，故
\[
r_{B}=8\sin\frac\pi3=4\sqrt3
\]
两点在同一射线上，所以
\[
|AB|=r_{B}-r_{A}=2\sqrt3
\]
\end{enumerate}
\end{solution}

\begin{problem}
设函数 \(f(x)=|x-a|+3x\)，其中 \(a>0\).
\begin{enumerate}
\item 当 \(a=1\) 时，求不等式 \(f(x)\geqslant 3x+2\) 的解集；
\item 若不等式 \(f(x)\leqslant 0\) 的解集为 \(\{x\mid x\leqslant -1\}\)，求 \(a\) 的值.
\end{enumerate}
\end{problem}

\begin{solution}
\begin{enumerate}
\item 当 \(a=1\) 时，分两段讨论：
\[
f(x)=|x-1|+3x=
\begin{cases}
1+2x,&x<1,\\
4x-1,&x\geqslant1.
\end{cases}
\]
故当 \(x<1\) 时，\(f(x)\geqslant3x+2\) 等价于 \(1+2x\geqslant3x+2\)，即 \(x\leqslant-1\)；当 \(x\geqslant1\) 时，\(4x-1\geqslant3x+2\)，即 \(x\geqslant3\). 解集为
\[
(-\infty,-1]\cup[3,+\infty)
\]
\item 对任意 \(a>0\)，当 \(x\geqslant a\) 时，
\[
f(x)=4x-a>0
\]
不满足 \(f(x)\leqslant0\)；当 \(x<a\) 时，
\[
f(x)=a+2x
\]
故 \(f(x)\leqslant0\) 等价于 \(x\leqslant-a/2\). 所以不等式解集为 \((-\infty,-a/2]\). 与题设的 \((-\infty,-1]\) 比较，得
\(-\frac a2=-1\)，\(a=2\)
\end{enumerate}
\end{solution}
